\documentclass[11pt]{article}
\usepackage[T1]{fontenc}
\usepackage{amsmath,amssymb,amsthm}
\usepackage{hyperref}
\usepackage[margin=1in]{geometry}

\newtheorem{theorem}{Theorem}
\newtheorem{corollary}{Corollary}
\newtheorem{lemma}{Lemma}

\newcommand{\R}{\mathbb R}
\newcommand{\E}{\mathbb E}
\newcommand{\Pp}{\mathbb P}
\newcommand{\Var}{\operatorname{Var}}
\newcommand{\supp}{\operatorname{supp}}

\title{A parity-separated partial result for anti-concentration under unconditional log-concavity}
\author{fa\_banach\_001, agent\_lane\_08}
\date{2026-07-18}

\begin{document}
\maketitle

\paragraph{Status.}
This is a \texttt{partial\_result\_likely\_valid} packet for arXiv:2603.22664.
It proves a structural subcase of Conjecture 1.4 in Glazer--Mikulincer, and
does not claim the full conjecture.

\paragraph{Source question.}
In the discussion after Theorem 1.3, Glazer--Mikulincer ask for uniform lower
bounds on the $L^2$ norm of degree-$d$ homogeneous polynomials under isotropic
log-concave measures invariant under signed coordinate permutations.  In their
Conjecture 1.4, with the implicit normalization $\operatorname{coeff}_d(f)=1$,
they predict a constant $C_d>0$ such that
\[
  \int_{\R^n} f^2\,d\mu \geq C_d
\]
for every isotropic log-concave $H_n$-invariant measure $\mu$ and every
degree-$d$ homogeneous polynomial $f$.  Question 1.5 asks whether the sharp
constant is governed, up to $O(1/n)$, by the isotropic cube.

\paragraph{Definitions.}
For a multi-index $I=(I_1,\ldots,I_n)$, write $x^I=\prod_j x_j^{I_j}$ and
$|I|=\sum_j I_j$.  A set $A$ of degree-$d$ multi-indices is
\emph{parity-separated} if for every distinct $I,J\in A$ there is a coordinate
$k$ such that $I_k+J_k$ is odd, equivalently $I\not\equiv J \pmod 2$
coordinatewise.

\paragraph{Proof intuition.}
The obstruction in the source conjecture is cancellation between monomials
that transform according to the same sign character.  If this cannot happen,
unconditionality does all the representation theory for free: after conditioning
on the absolute values $|X_i|$, the independent random signs kill every cross
term.  The problem reduces to proving that each fixed monomial has nontrivial
second moment.  That is a finite-dimensional consequence of log-concavity:
each one-dimensional isotropic log-concave marginal has uniformly bounded
density, so a fixed set of at most $d$ coordinates is simultaneously bounded
away from zero with probability bounded below in terms of $d$ only.

\begin{lemma}[Fixed monomial lower bound]
\label{lem:monomial}
There is a constant $m_d>0$ depending only on $d$ such that the following holds.
Let $X$ be an isotropic log-concave random vector in $\R^n$.  If $I$ is a
multi-index with $|I|=d$, then
\[
  \E |X^I|^2 \geq m_d .
\]
\end{lemma}

\begin{proof}
Let $S=\{j:I_j>0\}$, so $1\leq |S|\leq d$.  A standard one-dimensional fact says
that every isotropic log-concave real random variable has density bounded above
by a universal constant $L$.  Hence, for every coordinate $j$,
\[
  \Pp(|X_j|<a)\leq 2La .
\]
Choose $a=(4Ld)^{-1}$.  By the union bound,
\[
  \Pp\bigl(|X_j|\geq a\ \text{for all }j\in S\bigr)
  \geq 1-2L|S|a \geq \frac12 .
\]
On this event, $|X^I|^2=\prod_{j\in S}|X_j|^{2I_j}\geq a^{2d}$.  Therefore
$\E |X^I|^2\geq a^{2d}/2$, and the right-hand side depends only on $d$.
\end{proof}

\begin{theorem}[Parity-separated unconditional lower bound]
\label{thm:main}
For every $d\in\mathbb N$ there is a constant $c_d>0$ such that the following
holds.  Let $\mu$ be an isotropic unconditional log-concave probability measure
on $\R^n$.  Let
\[
  f(x)=\sum_{I\in A} a_I x^I
\]
be a degree-$d$ homogeneous polynomial whose support $A\subseteq\{I:|I|=d\}$ is
parity-separated.  Then
\[
  \int_{\R^n} f(x)^2\,d\mu(x)
  \geq c_d\sum_{I\in A}a_I^2 .
\]
In particular, if $\operatorname{coeff}_d(f)=1$, then
$\int f^2\,d\mu\geq c_d$.
\end{theorem}

\begin{proof}
Let $X\sim\mu$.  Since $\mu$ is unconditional, we may write
$X=\varepsilon\odot R$, where $R=(|X_1|,\ldots,|X_n|)$ and, conditionally on
$R$, the signs $\varepsilon_1,\ldots,\varepsilon_n$ are independent Rademacher
variables.  Then
\[
  \E\left[f(X)^2\mid R\right]
  =
  \sum_{I,J\in A} a_Ia_J R^{I+J}\,
       \E\left[\varepsilon^{I+J}\right].
\]
The sign expectation is zero unless $I\equiv J\pmod 2$.  Since $A$ is
parity-separated, the only surviving terms are the diagonal terms.  Thus
\[
  \E\left[f(X)^2\mid R\right]
  =
  \sum_{I\in A}a_I^2 R^{2I}.
\]
Taking expectations and applying Lemma \ref{lem:monomial} to each $I$ gives
\[
  \E f(X)^2
  =
  \sum_{I\in A}a_I^2\E |X^I|^2
  \geq
  m_d\sum_{I\in A}a_I^2 .
\]
Set $c_d=m_d$.
\end{proof}

\begin{corollary}[Multilinear subcase]
\label{cor:multilinear}
Conjecture 1.4 of Glazer--Mikulincer holds for all degree-$d$ multilinear
homogeneous polynomials under the weaker assumption that $\mu$ is isotropic,
unconditional, and log-concave.  If $f$ is multilinear and homogeneous with
$\operatorname{coeff}_d(f)=1$, then
\[
  \int f^2\,d\mu\geq c_d .
\]
Moreover, since every nonconstant multilinear monomial has an odd exponent in
some coordinate, $\E f=0$ under unconditional $\mu$.  Therefore
$\Var_\mu(f)\geq c_d$, and the equivalence Proposition 1.1 in the source paper
gives dimension-free small-ball and Fourier decay with exponent $1/d$.
\end{corollary}

\paragraph{Limitations.}
The theorem deliberately avoids the same-parity blocks.  In those blocks,
conditional sign averaging leaves cross terms
$a_Ia_J |X|^{I+J}$, so cancellations may persist.  This is precisely where the
source paper's higher-degree difficulty lives, especially for even symmetric
polynomials such as functions of $x_1^2,\ldots,x_n^2$.  The result also does
not address the cube-extremal constant in Question 1.5; the constant produced
above is only a robust dimension-free lower bound.

\paragraph{Novelty check.}
The local run indexes had no entry for arXiv:2603.22664.  A bounded web search
on 2026-07-18 for phrases including ``unconditional log-concave multilinear
polynomial small ball variance lower bound'', ``isotropic unconditional
log-concave homogeneous polynomial L2 norm coefficient lower bound'', and
``hyperoctahedral invariant log-concave homogeneous polynomial L2 norm
conjecture'' found the source paper and product-measure / $L^p$-ball results,
but no exact parity-separated unconditional diagonalization statement.

\paragraph{Human-review recommendation.}
Check the standard one-dimensional density bound used in Lemma
\ref{lem:monomial} and the conditional random-sign representation for
unconditional log-concave measures.  The rest of the proof is a direct
orthogonality computation.

\begin{thebibliography}{9}
\bibitem{GM}
I. Glazer and D. Mikulincer, \emph{Anti-concentration of polynomials:
$L^p$ balls and symmetric measures}, arXiv:2603.22664.

\bibitem{LV}
L. Lovasz and S. Vempala, \emph{The geometry of logconcave functions and
sampling algorithms}, Random Structures Algorithms 30 (2007), 307--358.
\end{thebibliography}

\end{document}
